%%%%%%%%%%  scan idx 368  |  folio 355  %%%%%%%%%%
Compte tenu de la définition de \(c\), et de (6.6.4), on obtient :
\[
  \widetilde c(\gamma^{i_j}(x_j))
  =
  \left(0,1+(-1)^{i_j-1}(i_j-1)!\,\gamma^{i_j}(x_j)+\cdots\right)
\]
D'après (6.5.3), et le fait que \(\widetilde c(a)=(a,1)\), on obtient finalement :
\[
\begin{aligned}
\widetilde c(z)
&=\sum
\left(0,\left(1+
(-1)^{i_1+\cdots+i_k-1}
(i_1+\cdots+i_k-1)!
\,\gamma^{i_1}(x_1)\cdots\gamma^{i_k}(x_k)
+\cdots\right)^a\right)\\
&=
\left(0,1+
\sum_{\substack{(x_1,\ldots,x_k)\\ i_1+\cdots+i_k=m}}
(-1)^{m-1}(m-1)!a\,\gamma^{i_1}(x_1)\cdots\gamma^{i_k}(x_k)
+\cdots\right)\\
&=
(0,1+(-1)^{m-1}(m-1)!\bar z+\cdots),
\end{aligned}
\]
d'où le résultat.

\textbf{Proposition 6.10.} Soient \(K\) un anneau binômial, \(R\) une \(K\)-\(\lambda\)-algèbre augmentée, \(N\) un élément de \(R\) tel que \(\lambda^k(N)=0\) pour \(k>d\), et \(x\in\Fil^iR\). Alors :
\[
  \gamma^{d+i}(N,x)-(-1)^{d+i-1}(d+i-1)!x\in\Fil^{i+1}R.
\]

On peut trouver des éléments \(x_{\alpha,j}\in R\), et \(a_\alpha\in K\) tels que :
\[
  x=\sum_\alpha a_\alpha\,
       \gamma^{i_1}(x_{\alpha,1})\cdots
       \gamma^{i_k}(x_{\alpha,k}),
\]
avec \(i_1+\cdots+i_k\ge i\), et \(\epsilon(x_{\alpha,j})=0\) pour tout \(\alpha\) et tout \(j\). Soit \(R_0\) la \(K\)-\(\lambda\)-algèbre engendrée par des générateurs \(N'_0\), \(X_{\alpha,j}\) soumis aux relations :
\[
  \gamma^k(N'_0)=0\quad\text{pour }k>d.
\]
On définit un homomorphisme d'augmentation \(R_0\to K\) par \(\epsilon(N'_0)=\epsilon(X_{\alpha,j})=0\). D'après (4.9.1), \(R_0\) est l'anneau des polynômes à coefficients dans \(K\) par rapport aux indéterminées \(\gamma^n(N'_0)\), avec \(1\le n\le d\), et \(\gamma^p(X_{\alpha,j})\) avec \(p\ge1\), et la \(\lambda\)-filtration de \(R_0\) est identique à sa filtration par le poids, en attribuant le poids \(k\) à \(\gamma^k(N'_0)\) et \(\gamma^k(X_{\alpha,j})\) (4.12 et 4.15).

On pose :
\[
  X_0=\sum_\alpha a_\alpha\,
       \gamma^{i_1}(X_{\alpha,1})\cdots
       \gamma^{i_k}(X_{\alpha,k}),
\]
